\documentclass[a4paper,addpoints,12pt]{exam}
%\documentclass[a4paper,answers,addpoints,12pt]{exam}
\unframedsolutions

%entorno para química: texto y math mode
%\newcommand*\chem[1]{\textup{#1}}
\newcommand*\chem[1]{\ensuremath{\mathrm{#1}}}

%tamaños fuentes extra: 8,9,14,17,20
	\usepackage{extsizes}

\usepackage[dvipsnames]{xcolor}
\usepackage{array,amsmath,fancybox,amsfonts,graphicx,amssymb,wasysym,latexsym,slashed,subfigure,calc}
%símbolo bonito para qed
\usepackage{amsthm}
	%\newcommand{\mc}[3]{\multicolumn{#1}{#2}{#3}}
\newcommand{\halmos}{\vspace*{-4mm}
	
					\begin{flushright}
						\rule{1mm}{2.5mm} 
					\end{flushright}\vspace*{-4mm}
						
					  }
\renewcommand{\qed}{\halmos}

\usepackage{mdframed}
	
%caja bonita
  \usepackage[many]{tcolorbox}
\definecolor{titlebg}{RGB}{80,100,72}
\newtcolorbox{learn}{
  breakable,
  enhanced,
  arc=0pt,
  outer arc=0pt,
  colframe=titlebg,
  colback=titlebg!7,
  overlay unbroken and first={
    \node[
      draw=titlebg,
      fill=titlebg,
      rotate=270,
      anchor=north west,
      text=white,
      font=\bfseries
    ]
    at (frame.north west)  
    {DATOS};
  }
}

\usepackage{anysize}
	%\marginsize{left}{right}{top}{bottom}
    \marginsize{2cm}{2cm}{1.5cm}{2cm}
    
\usepackage[spanish]{babel}
\usepackage[utf8]{inputenc}

\usepackage{sidecap}
\usepackage{mwe,float}

%redefiniciones
\pointpoints{punto}{puntos}
\bonuspointpoints{punto}{puntos}
\hqword{Pregunta} %substitutes text for \Question"
%\vpgword{Pág.}%substitutestext for \Page"
\hpword{Máx.}%substitutes text for \Points"
\hsword{Nota}%substitutes text for \Score"
\htword{Total}%substitutes text for \Total:"

\chqword{Pregunta}
\chpword{Máx.}
\chsword{Nota}
\chtword{Total}
\chbpword{Puntos \textit{bonus}}

\renewcommand{\solutiontitle}{\noindent\textbf{Solución:}\par\noindent}



%versores: %uso: \uveci\ne\hat{\imath}_{\uvec{i}}+\uvec{j}+\uvec{k}$
	\usepackage{bm}
	\newcommand{\uveci}{{\bm{\hat{\textnormal{\bfseries\i}}}}}
	\newcommand{\uvecj}{{\bm{\hat{\textnormal{\bfseries\j}}}}}
	\DeclareRobustCommand{\uvec}[1]{{%
 	 \ifcsname uvec#1\endcsname
  	   \csname uvec#1\endcsname
  	 \else
   		 \bm{\hat{\mathbf{#1}}}%
   		\fi
									}}

%fuentes
%	\usepackage[T1]{fontenc}
%	\usepackage{fourier}


\usepackage[T1]{fontenc}
\usepackage{mathpazo}

	%\usepackage[T1]{fontenc}
	\usepackage{librebaskerville}

\usepackage{amssymb,amsmath}
\usepackage{graphicx,graphics}
\usepackage{color}
	%	\colorgrids %grids con color
	%	\definecolor{GridColor}{gray}{0.8}
		%grid
%			\setlength{\gridsize}{5mm}
%   			 \setlength{\gridlinewidth}{0.1pt}

%fracciones en texto más elaboradas
	\usepackage{xfrac}

%mejora en la exposición de guarismos
	\usepackage{numprint}
 
%caption
\usepackage[font=small,labelfont=bf,
labelsep=endash]{caption}
	
%soluciones sombreadas	
\shadedsolutions


\checkboxchar{$\Box$}
\checkedchar{$\blacksquare$}

\usepackage{everypage,draftwatermark} %Marca de agua
\SetWatermarkText{\parbox{2\textwidth}{\centering IES\\Gerardo\\ Diego}}
\SetWatermarkScale{1.0}
\SetWatermarkLightness{0.95}



\pagestyle{headandfoot}
%\runningheadrule
%\firstpageheader{Física para Biólogos}{Boletín 1a}{1 de octubre de 2011}
%\runningheader{Física para Biólogos}
%              {Boletín 1a, Página \thepage\ de \numpages}
%            {1 de octubre de 2011}
%\firstpagefooter{}{}{}
%\runningfooter{}{}{}


\header{%\includegraphics[scale=.25]{iesgd2.jpeg} 
		\hspace*{.5cm}\vspace*{-.4cm}
		\includegraphics[scale=.35]{a.jpeg}}       
      {%\\Noviembre de 2019
       }
       {\large \textit{Global}\textbf{B2}}
       %aleph, beth, gimel, daleth
       
%pie de página

   %enumera en todas menos en la primera
   %\runningfootrule

  % enumera en todas
  % \footrule
   \footer{}{}{}
   %\footer{}{Página \thepage\ de \numpages}{\iflastpage{}}
   
%página blanco
\usepackage{afterpage}
%uso \afterpage{\blankpage}

\newcommand\blankpage{%
    \null
    \thispagestyle{empty}%
    \addtocounter{page}{0}%
    \newpage}


\begin{document}
%\pagenumbering{gobble}
%formato de preguntas
\qformat{\textbf{Pregunta \thequestion}\quad (\thepoints)\hfill}

\thispagestyle{empty}
\vspace*{-2cm}

 {\begin{center} \Huge FÍSICA---\textit{Global}\textbf{B2}\end{center}}
 \vspace{1cm}
 
\fbox{\parbox{17cm}{\textbf{Nombre y apellidos}:
\vspace*{10mm}
				}
	}
\vspace{1mm}

\fbox{\parbox{10cm}{\textbf{Grupo, subgrupo}:
\vspace*{10mm}
				}
	}


%\boxed{\textit{\hspace*{5mm} Sólo una de las respuestas es verdadera. Los errores restan 0.25 puntos.}}

\vspace*{\fill}
	
\begin{learn}
{
\vspace{5mm}
$$\triangleright\;h\simeq 6.626\cdot 10^{-34}\;J\,s\hspace{1.1cm} \triangleright\;c_0\simeq 3\cdot 10^8\;m\,s^{-1}\hspace{1.1cm} \triangleright\;\hbar c\simeq 197\;eV\,nm$$
\vspace{-5mm}

$$\triangleright\;m_{e}c_0^2\simeq 511\;keV\hspace{1cm} 1\;eV\simeq 1.6\cdot 10^{-19}\;J$$
\vspace{-5mm}

$$\triangleright\;K\simeq 9\cdot 10^{9}\; N\, m^2\, C^{-2}\hspace{1cm}\triangleright\;G\simeq 6.67\cdot 10^{-11}\;N\,m^2\,kg^{-2}$$
\vspace{-7mm}

$$\triangleright\; M_{\text{Júpiter}}\simeq 1.9\cdot 10^{27}\;kg\hspace{1cm}
\triangleright\;R_{\text{Júpiter}}\simeq 69\, 911\;km\hspace{1cm} \triangleright\;m_{\text{Juno}}\simeq 3\,625\;kg	$$			
 }
\end{learn}



\newpage

\begin{questions}
	\begin{question}[2]
		\textbf{(Madrid, Ordinaria, 2011)} Considera un material conductor de forma esférica y radio $R=10\;cm$, cargado con una densidad superficial de carga homogénea igual a $\sigma=+1.18\cdot 10^4\;nC\,m^{-2}$.
		\begin{parts}
			\part Calcula el \textbf{vector campo electrostático} creado en los puntos situados a distancias de $5\;cm$ y $10\;cm$ del centro de la esfera.
			\begin{mdframed}[backgroundcolor=brown!10] 
			El campo en el interior de un conductor es nulo. El campo en la superficie puede calcularse con el teorema de Gauss: el valor es el mismo que se encontraría concentrando toda la carga de la superficie en el centro de la esfera, \textit{i.e.}
			$$\vec{E}\big|_{r=R}=\frac{\sigma}{\epsilon_0}\uvec{u}_r=4\pi K\;\sigma\;\uvec{u}_r,$$
			donde $\uvec{u}_r$ es un vector radial dirigido al exterior de la esfera.
			\end{mdframed}
			\part ¿Cuál es el \textbf{trabajo} involucrado en el desplazamiento de una carga de $2\;nC$ desde una distancia muy grande a la superficie de la esfera?
			\begin{mdframed}[backgroundcolor=brown!10] 
			Si la distancia es muy grande comparada con las dimensiones de la esfera podemos suponer que $r_{inicial}\to\infty$, de manera que en el proceso de desplazamiento el trabajo involucrado será
			$$W=-\Delta E_p=-E_{p,f}=-\frac{q_{ext}\;4\pi K\;R^2\sigma}{R}=-4\pi K\; R\,\sigma \,q_{ext}<0,$$
			como corresponde a un proceso evidentemente no espontáneo.
			\end{mdframed}
					
		\end{parts}
	\qed
	\end{question}
	
	\newpage
	\begin{question}[2]
	\textbf{(Cataluña, Extraordinaria, 2019)} Desde el día 5 de julio de 2016 la sonda \textit{Juno} describe una órbita polar alrededor de \textit{Júpiter} con objeto de estudiar características atmosféricas y estructurales del gigante gaseoso. Suponiendo que la órbita sea circular y sabiendo que su altura sobre el planeta es de $4\,300\;km$:
		\begin{parts}
			\part Calcula la \textbf{energía cinética} de \textit{Juno} y su periodo de rotación.
			\begin{mdframed}[backgroundcolor=brown!10] 
			A partir del teorema del virial para órbitas ligadas circulares es posible escribir
			$$E_c=-\frac{E_p}{2}=\frac{1}{2}G\frac{Mm}{R+h}\;,$$
			con $R$ el radio de Júpiter y $h$ la altura del satélite sobre la superficie.
			Por otro lado se cumplirá la ley de Kepler, luego
			$$T=\sqrt{\frac{4\pi^2(R+h)^3}{GM}}.$$
			\end{mdframed}
			\part Calcula la \textbf{energía} que habría que comunicar a \textit{Juno} para que abandone el campo gravitatorio de \textit{Júpiter}.
			\begin{mdframed}[backgroundcolor=brown!10] 
			La energía necesaria será, prescindiendo de convenios para el signo, la energía mecánica que tiene en la órbita: es precisamente esto lo que debemos hacer para llevarlo al infinito con energía mínima.
			Con todo, y aplicando el teorema del virial, obtenemos
			$$W=|E_{m,\text{ órbita}}|=\frac{E_p}{2}=\frac{1}{2}G\frac{Mm}{R+h}\;,$$
			que es la mitad, exactamente, de la energía mínima que deberíamos aplicarle si permaneciera estático en esa posición\footnote{Recuerda que \textbf{ya} tiene velocidad en la órbita, luego la velocidad a la que tendríamos que impulsarle es menor que la velocidad de escape correspondiente al satélite estático en ese mismo punto.}.
			\vspace{1cm}
			
			\end{mdframed}				
		\end{parts}
		
	\qed
	\end{question}
	
	\newpage
	\begin{question}[2]
	\begin{minipage}{8.5cm}
			\textbf{(Cataluña, Ordinaria 2019)} Se sitúa una partícula puntual con una carga eléctrica de $q=1.9\;\mu C$ y dos partículas puntuales iguales y de carga $-q$ en los vértices de un cuadrado de lado $a=1.50\;cm$ tal como se indica en la figura. 
	\end{minipage}\hfill
	\begin{minipage}{6cm}
		\vspace{-7mm}
		
		\def\svgwidth{4.9cm}
			\input{cuadrado.pdf_tex}
	\end{minipage}
	\vspace{2mm}
	
		\begin{parts}
			\part ¿Cuál es el\textbf{ valor de la carga} $q_0$ situada en el vértice libre si la fuerza sobre la partícula de carga $q$ es nula?
			\begin{mdframed}[backgroundcolor=brown!10] 
			Colocamos el sistema de referencia con origen en la carga $q$ y orientado positivamente \textit{hacia arriba} y \textit{hacia la derecha.} 
			
			De esta manera, para imponer la anulación de la fuerza sobre $q$ debido al resto de cargas podemos invocar la superposición lineal de campos: 
			$\vec{F}(q)=\vec{F}_{-q}+\vec{F}_{-q}+\vec{F}_{q_0}=0.$
			Desarrollando vemos que ha de ser $K\frac{q^2}{a^2}(-\uvec{i})+K\frac{q^2}{a^2}(-\uvec{j})+\vec{F}\big|_{q_0}=0,$
			o lo que es lo mismo
			$\vec{F}_{q_0}=\frac{q^2}{a^2}(\uvec{i}+\uvec{j}).$
			Escribiendo, por otro lado, esta misma fuerza en términos de la fuente $q_0$ obtenemos
			$\vec{F}_{q_0}=K\frac{qq_0}{(\sqrt{2}a)^2}\frac{1}{\sqrt{2}}(\uvec{i}+\uvec{j}).$
			Comparando con la expresión anterior llegamos a $q_0=2\sqrt{2}q$.
			\end{mdframed}
			\part Calcula el \textbf{trabajo} involucrado al cambiar de posición una carga puntual de $0.50\;\mu C$ desde un punto muy lejano hasta el centro del cuadrado.
			\begin{mdframed}[backgroundcolor=brown!10] 
			Si consideramos que la distancia inicial es mucho más grande que las dimensiones de la configuración de cargas es lícito tomar $r_{inicial}\to\infty$. El trabajo se calculará superponiendo todas las contribuciones 
			\begin{eqnarray*}
			W=-\Delta E_{p\,total}=0-E_{p\,total}^{(f)}\!\!=-\!\!\sum_{\substack{las\,cuatro\\cargas}}\!\!E_{p}^{(f)}&=&-K\,q_{ext}\left(-2\frac{q}{\frac{a}{\sqrt{2}}}+\frac{q_0}{\frac{a}{\sqrt{2}}}+\frac{q}{\frac{a}{\sqrt{2}}}\right)\\&=&-\sqrt{2}K\frac{q_{ext}}{a}(q_0-q)\\&=&-\frac{7\sqrt{2}}{1+2\sqrt{2}}\;K\frac{q\,q_{ext}}{a},
			\end{eqnarray*}
			donde $\sfrac{a}{\sqrt{2}}$ es la longitud de la mitad de la diagonal del cuadrado\footnote{El coeficiente irracional es un factor geométrico que depende de la estructura de la configuración, mientras que el término proporcional a las cargas e inversamente proporcional a la longitud característica de la configuración hace las veces de\textit{ escala típica} para la energía potencial electrostática sistema: cualquier energía involucrada debe ser un múltiplo de ella.}.
			\vspace{1cm}
			
			\end{mdframed}
		\end{parts}
	\qed
	
	\end{question}
	
	\newpage
	\begin{question}[2]
		
		\begin{parts}
			\part \textbf{(Cataluña, Ordinaria, 2016)}. La irradiancia solar que llega a la superficie de la Tierra (es decir: \textit{la potencia incidente por unidad de superficie}) es aproximadamente $ 1400\;W\,m^{-2}$. Suponiendo que la energía media de los fotones que llegan es de $2.2\;eV$, calcula la \textbf{longitud de onda media} de los fotones y el \textbf{número} de ellos que incide sobre una superficie de $1.00\;cm^2$ cada segundo.
			\vspace{3mm}
			\begin{mdframed}[backgroundcolor=brown!10] 
			La longitud de onda media de los fotones será, en virtud de la ley de Planck, 
			$$E=\frac{hc}{\lambda}\rightarrow \lambda=\frac{hc}{E}\simeq\frac{1240}{2.2}\simeq 564\,nm.$$
			En un segundo y $1\;cm^2$ la energía total depositada por los fotones es
			$E=1400\cdot 10^{-4}\;J.$ Como cada fotón tiene una energía de $2.2\cdot 1.6\cdot 10^{-19}\;J$ el número de fotones en esas condiciones es 
			$$N\equiv\frac{0.14}{2.2\cdot 1.6\cdot 10^{-19}}\simeq 4\cdot 10^{17}.$$
			
			\end{mdframed}
			
			\part \textbf{(Madrid, Ordinaria, 2015)} Cuando se encuentra fuera del núcleo atómico, el neutrón es una partícula inestable con una vida media aproximada de $885.7\;s$. Una fuente de neutrones emite $10^{10}$ neutrones por segundo con una velocidad constante de $100\;km\,s^{-1}$,  ¿qué \textbf{proporción} de los neutrones emitidos en un segundo recorren, en promedio, una distancia de $4\cdot 10^5\;km$ \textit{sin desintegrarse}?
			\begin{mdframed}[backgroundcolor=brown!10] 
			¿Cuánto tarda el haz en recorrer $\Delta x=4\cdot 10^5\;km$? Al ser un movimiento uniforme el cálculo es sencillo: $t=\sfrac{\Delta x}{v}=4\cdot 10^5\cdot 10^{-2}=4\cdot 10^3\;s.$ Ahora debemos ver qué fracción de neutrones queda cuando ha transcurrido ese tiempo: a partir de la ley de desintegración nuclear, y asumiendo una probabilidad por unidad de tiempo constante, resulta
			$$\frac{N(t)}{N_0}=\exp\left(-\frac{t}{\tau}\right)=\exp\left(-\frac{\Delta x}{v\,\tau}\right),$$
			con $\tau$ la vida media del neutrón fuera del núcleo.
			\end{mdframed}
					
		\end{parts}
	\qed
	\end{question}
	
	\newpage
	
	\begin{question}[2]
		\textbf{(Islas Canarias, Ordinaria, 2017)} Una lente convergente de un proyector de diapositivas que tiene una distancia focal de $16\;cm$ proyecta nítidamente la imagen de una diapositiva de $3\;cm$ de alto sobre una pantalla que se encuentra a $4\;m$ de la lente.
		\begin{parts}
			\part ¿A qué \textbf{distancia} de la lente está colocada la diapositiva y cuál es el \textbf{aumento} de la imagen formada por el proyector en la pantalla?
			\begin{mdframed}[backgroundcolor=brown!10] 
			Aplicando la ecuación de las lentes delgadas con $f'=+16\;cm$, como corresponde a una lente convergente, y $s'=+400\;cm$ resulta
			$$\frac{1}{f'}=\frac{1}{s'}-\frac{1}{s}\longrightarrow s=\left(\frac{1}{s'}-\frac{1}{f'}\right)^{-1}.$$
			El aumento, una vez hallado este dato, será
			$M=\sfrac{s'}{s},$
			que, con seguridad, será negativo debido a que la imagen formada es real, según se interpreta del problema.
			\end{mdframed}
			\part Realiza el \textbf{diagrama de rayos} de la situación planteada.
				\begin{mdframed}[backgroundcolor=brown!10] 
				Corresponde a la construcción para una lente convergente a partir de un objeto real situado a la izquierda (según convenio usado en clase) del foco objeto.
				\end{mdframed}
		\end{parts}
	\qed
	\end{question}
	
	\newpage
	\begin{question}[2]
	\textbf{(Cantabria, Ordinaria, 2013)}. Por una cuerda muy larga se propaga un movimiento ondulatorio caracterizado por la ecuación:
	$$y(x,t)=2\sin\left(2\pi\left[\frac{t}{4}-\frac{x}{2}\right]\right),\hspace{.5cm}\text{\small(SI).}$$
		\begin{parts}
			\part Halla el \textbf{periodo}, la \textbf{frecuencia}, el \textbf{número de onda} y la \textbf{velocidad de propagación} de la onda armónica.
			\begin{mdframed}[backgroundcolor=brown!10] 
			A partir de la onda provista se deduce, por inspección directa, 
			$$\omega=\frac{\pi}{2}\;s^{-1}\hspace{.3cm} k=\pi\,m^{-1}\hspace{.3cm},$$
			de manera que $T=\sfrac{2\pi}{\omega}=4\;s$ y $c=\sfrac{\omega}{k}=0.5\;m\,s^{-1}.$
			\end{mdframed}
			\part Halla la \textbf{distancia mínima} a la que se encuentran, en un instante fijo cualquiera, dos puntos de la cuerda cuya diferencia de fase es $\sfrac{3\pi}{2}$ radianes.
			\begin{mdframed}[backgroundcolor=brown!10] 
			La fase de la onda será $\delta\equiv |\omega \Delta t-k\Delta x|.$ Para $\Delta t=0$ resulta que $\delta=k\,\Delta x$. Asumiendo que el desfase prescrito se encuentra en la \textit{primera vuelta} la distancia será mínima e igual a 
			$$\Delta x=\frac{\delta}{k}.$$
			\end{mdframed}
					
		\end{parts}
	\qed
	\end{question}
	
	\newpage
	\begin{question}[2]
		\textbf{(Cataluña, Extraordinaria, 2017)} En un estadio de fútbol, el público hace \textit{la ola} para celebrar la buena actuación del equipo local. Los espectadores vuelven a levantar los brazos cada $10\;s$ y \textit{la ola} producida es tan grande que cualesquiera dos de ellos sentados en la misma fila y separados por una distancia de $50\;m$ se mueven sincrónicamente. Modelizando \textit{la ola} como una onda armónica:
		\begin{parts}
			\part  ¿Qué \textbf{tipo de polarización} presenta? Calcula su \textbf{longitud de onda}.			
			\begin{mdframed}[backgroundcolor=brown!10] 
			\textit{La ola} se forma al mover los brazos de arriba a abajo, de manera que la vibración de las partículas es ortogonal a la dirección de propagación: se trata de una onda transversalmente polarizada. 
			
			La longitud de onda es precisamente la distancia horizontal mínima que conecta dos puntos en fase: \textit{i.e.} $\lambda=50\;m$.
			\end{mdframed}
			\part Las manos de cada espectador, al pasar \textit{la ola} y levantarse, se mueven verticalmente $1.0\;m$ desde el punto más bajo al más alto. Escribe la \textbf{ecuación de la onda} que se propaga en una misma fila asumiendo que $y(0,0)=0$.
			\begin{mdframed}[backgroundcolor=brown!10] 
			La amplitud será \textit{la mitad} de la distancia recorrida entre las posiciones más alta y más baja. Por lo demás, tomando el desfase inicial como nulo, tenemos
			$$y(x,t)=0.5\sin\left(kx\pm\omega t\right),(SI),$$
			donde $k=\sfrac{2\pi}{\lambda}$ y $\omega=\frac{2\pi}{T}$, siendo $T=10\;s$, pues este es precisamente el tiempo en el que las manos de un espectador cualquiera vuelven a estar en el mismo estado de vibración.
			\end{mdframed}
					
		\end{parts}
	\qed
	\end{question}
	
	\newpage
	\begin{question}[2]
		
		\begin{parts}
			\part \textbf{(Cataluña, Ordinaria, 2018)} El polonio, $\chem{^{210}_{\;\;84}Po}$, es un emisor natural de partículas $\alpha$. Sabiendo que su periodo de semidesintegración es de $138$ días, calcula la \textbf{masa de polonio} que queda en una muestra de $10.0\;g$ después de $69$ días desde el inicio de la toma de medidas.
						\vspace{3mm}
			
			\begin{mdframed}[backgroundcolor=brown!10] 
			La ley de desintegración de una partícula inestable que tiene la misma probabilidad de desintegrarse por unidad de tiempo nos permite escribir la expresión
			$$m(t)=m_0\exp\left(\frac{\ln(2)}{T}\,t\right),$$
			para la masa de polonio no desintegrado transcurrido un tiempo $t$.
			\end{mdframed}
			\part  \textbf{(Cataluña, Ordinaria, 2009)} Una fuente luminosa emite luz monocromática de longitud $550\;nm$, con una potencia de $2\;mW$. Se hace incidir la luz sobre un metal produciendo la emisión de electrones. La energía mínima de extracción de los electrones de la red metálica es de $2.10\;eV$. Calcula la \textbf{velocidad} del electrón más rápido que se emite y el \textbf{potencial de detención} necesario para anular la corriente fotoelectrónica.
			\begin{mdframed}[backgroundcolor=brown!10] 
			El balance energético del efecto fotoeléctrico puede escribirse como
			$E_{incidente}=E_{c,max}+W,$
			de manera que la energía cinética máxima será $$E_{c,max}=\frac{hc}{\lambda}-W\simeq \frac{1240}{550}-2.1=0.16\;eV.$$
			Escribiendo explícitamente la energía cinética en función de la velocidad y asumiendo 				que el flujo fotoelectrónico es no-relativista, resulta
			$$v=\sqrt{\frac{2E_{c,max}}{m_e}}=\sqrt{\frac{2E_{c,max}c^2}{m_ec^2}}\simeq \sqrt{\frac{2\cdot 0.16\cdot(3\cdot 10^8)^2}{5.11\cdot 10^5}}\simeq 2.3\cdot 10^5\;m\,s^{-1},$$
			una vez hallada la masa del electrón con los datos provistos.
			
			Naturalmente, para detener al fotoelectrón más veloz el potencial de detención cumplirá
			$$e\Delta V_{stopping}=E_{c,max}\longrightarrow \Delta V_{stopping}\simeq\frac{0.16\;eV}{e}=0.16\;V.$$
			\end{mdframed}
					
		\end{parts}
	\qed
	\end{question}
	
	\afterpage{\blankpage}
	
\end{questions}


\end{document}






